\subsection{Résolvantes du transport et conservation de la mémoire}
\label{subsec:ii-memoria-resolvente}

L’histoire prolongée détermine les fenêtres
\(I_m=\{9m+1,\ldots,9m+9\}\), pour \(m\geq0\).
Avec la numérotation de cette section, \(I_m=E_m\) pendant le premier cycle.
Soit \(\sigma_m\in\{-1,1\}\) l’orientation conservée par cette histoire ;
pour \(0\leq m<12\), elle coïncide avec \(\Sigma_{12}(m)\). On définit
\[
 a_m=\sigma_m\sum_{t\in I_m}\mathbf1_{\{2,4,6,8\}}(d_t),
 \qquad |a_m|\leq9.
\]
Le registre et le repère mobile se prolongent par
\[
 \mathbf z_{m+1}=\mathbf z_m+a_m\mathbf e_{m\bmod12},\qquad y_m=S^{-m}\mathbf z_m.
\]
L’état initial \(\mathbf z_0\) conserve l’histoire antérieure ; il vaut zéro pour
un préfixe initialement vide. Dans cette partie, nous écrivons
\(R=S^{-1}\) pour le transport du repère. La mise à jour
\eqref{eq:eta} est maintenue dans toute fenêtre. L’orientation et les événements
ultérieurs s’obtiennent à partir de l’histoire prolongée : le retour du repère
ne présuppose pas la périodicité de la suite d’événements.

\begin{proposition}[Identité de troncature avec état terminal]
Pour $N\geq1$, soient
\[
Y_N(u)=\sum_{m=0}^N y_mu^m,\qquad
H_{\eta,N}(u)=\sum_{m=0}^{N-1}\eta_mu^m.
\]
On a l’identité polynomiale exacte
\begin{equation}
(I-uR)Y_N(u)=y_0+uH_{\eta,N}(u)-u^{N+1}Ry_N.
\label{mem:identidad-finita}
\end{equation}
\end{proposition}
\begin{proof}
Aux degrés $1\leq m\leq N$, le coefficient du membre de gauche est
$y_m-Ry_{m-1}=\eta_{m-1}$, d’après \eqref{eq:eta}.
Ses coefficients aux degrés zéro et $N+1$ sont $y_0$ et $-Ry_N$,
respectivement. La comparaison des coefficients prouve l’égalité.
\end{proof}

Le terme terminal conserve la mémoire transportée lorsqu’on coupe une histoire.
L’identité permet de comparer une évaluation finie à son prolongement
en maintenant explicite l’état qui reste à la frontière.

\subsubsection{Correspondance avec le transport du noyau commun}

La numérotation des fenêtres de la section~\ref{subsec:memoria-resolvente}
et celle employée ici décrivent la même partition de l’histoire. La
correspondance est
\[
 \underbrace{I_m=E_{m+1}}_{\text{indice du noyau commun}}
 \quad\longleftrightarrow\quad
 \underbrace{I_m=E_m}_{\text{indice de cette section}},
 \qquad 0\leq m<12.
\]
L’indice change le nom de la fenêtre ; ses neuf transitions,
son orientation \(\sigma_m=\Sigma_{12}(m)\), son événement \(a_m\) et
sa position dans le registre restent identiques. Pour tout prolongement,
\(R=S^{-1}\), \(y_m=S^{-m}\mathbf z_m\) et la
mise à jour~\eqref{eq:eta} coïncident avec
\eqref{mem:actualizacion}. L’orientation et la suite d’événements
conservent la dépendance à l’égard de l’histoire, également après le
retour du repère.

Cette identification permet d’appliquer directement la
proposition~\ref{mem:prop-serie} : la série~\eqref{mem:serie}, la
borne~\eqref{mem:cota} et la convergence uniforme de ses dérivées
correspondent à ce même registre et à la même horloge de fenêtres. L’identité
finie~\eqref{mem:identidad-finita} détermine en outre le terme
terminal lorsqu’on tronque cette série. Ses hypothèses maintiennent l’histoire initiale
et l’événement borné \(|a_m|\leq9\).

En réutilisant le registre, la proposition~\ref{mem:prop-schur} conserve
conjointement l’opérateur, la source et le lecteur au moyen de
\eqref{mem:schur-fuente} et~\eqref{mem:schur-lector}. L’exemple
\eqref{mem:ejemplo-ciclo} précise l’effet des excursions omises
par une compression. La composition affine~\eqref{mem:cociclo-tres}, les
blocs intermédiaires~\eqref{mem:schur-intermedio} et l’identité
\eqref{mem:schur-transitivo} garantissent la cohérence lors de la concaténation de
segments ou de l’élimination de plusieurs couches. Leurs démonstrations complètes figurent
dans la section~\ref{subsec:memoria-resolvente}, avec les domaines, les
sources et les applications de lecture qui y sont spécifiés.

Le bilan~\eqref{ii:mem:balance} conserve la correspondance entre les
lecteurs orientés du registre signé et la forme quadratique commune.
La distinction entre le vecteur accumulé \(\mathbf z_m\) et le bloc
décimal d’incidence \(z_m\) reste celle établie dans cette section.
Les identités de transport s’appliquent à la mémoire ainsi typée ; elles ne
déterminent pas à elles seules des coefficients analytiques supplémentaires.
